\documentclass[11pt,a4paper]{article}
\usepackage[margin=1in]{geometry}
\usepackage{amsmath}
\usepackage{graphicx}
\usepackage{hyperref}
\usepackage{cite}

\title{A Two-Stage Search Engine for Python Code using BM25 Retrieval and Dense Re-ranking}
\author{Tatyana Shmykova \and Albina Akhmetova \and Dmitrii Naumov}
\date{\today}

\begin{document}
\maketitle

\section{Problem Statement and Motivation}

Developers often search for code using natural language queries such as ``read csv file with pandas'' or ``sort dictionary by value''. Standard keyword search does not always work well for this task, because the query words and the code tokens may be different. For example, a query may use the phrase ``remove duplicates'', while the code implementation may use different variable names or functions.

The goal of this project is to build a small search engine for Python functions that retrieves relevant code snippets from a natural language query. The system will use a two-stage retrieval pipeline. In the first stage, a classical lexical retrieval method BM25 will retrieve a large candidate set of documents. In the second stage, a dense semantic model will compute similarity between the query and the candidate documents and rerank them.

This approach allows combining the efficiency of lexical retrieval with the semantic understanding of dense embeddings. The project will investigate whether semantic reranking improves the quality of search results compared to a lexical baseline.

\section{Dataset and Corpus}

The dataset used in this project will be the Python subset of the CodeSearchNet dataset. CodeSearchNet is a large benchmark dataset designed for natural-language-to-code retrieval tasks. It contains code functions extracted from open-source repositories together with their documentation and natural language descriptions.

In this project, each Python function will be treated as a document. The textual representation of a document will include:
\begin{itemize}
    \item the function name,
    \item the docstring (when available),
    \item the code body.
\end{itemize}

The natural language queries provided in the dataset will be used for evaluation. The dataset also contains relevance labels that indicate which functions correspond to a given query.

To keep the system manageable, the project will focus only on the Python portion of the dataset. If necessary, a smaller subset of the data will be used during development.

\section{Proposed Methodology}

The search engine will follow a two-stage retrieval architecture.

\subsection{Stage 1: BM25 Retrieval}

In the first stage, the system will build an inverted index over the corpus of Python functions. Each function will be tokenized and indexed using standard information retrieval preprocessing steps such as tokenization and normalization.

The BM25 ranking function will be implemented and used to retrieve the top $K$ candidate documents for a given query.

The goal of this stage is to provide a high-recall candidate set that contains potentially relevant code snippets.

\subsection{Stage 2: Dense Re-ranking}

In the second stage, the system will rerank the candidate documents returned by BM25 using dense semantic representations.

For this step, the project will use the pretrained CodeBERT model \cite{codebert}. CodeBERT is a transformer model trained jointly on natural language and source code data, including the CodeSearchNet dataset. The model produces vector embeddings for both queries and code snippets in a shared semantic space.

Both the natural language query and the candidate Python functions will be encoded using CodeBERT. Similarity between the query embedding and document embeddings will be computed using cosine similarity.

Only the $K$ candidates retrieved by BM25 will be encoded and scored in this stage. The documents will then be reranked according to their dense similarity scores, and the final top $k$ results will be returned.

Using dense reranking instead of full dense retrieval significantly reduces computational cost while still allowing semantic matching between queries and code snippets.

\section{Implementation Plan}

The system will be implemented in Python. The main components will include:

\begin{itemize}
    \item parsing and preprocessing the CodeSearchNet Python dataset,
    \item building an inverted index for BM25 retrieval,
    \item implementing the BM25 ranking function,
    \item generating dense embeddings for queries and documents using the CodeBERT model,
    \item reranking candidate documents based on cosine similarity between embeddings.
\end{itemize}

Standard Python libraries such as NumPy, scikit-learn, and PyTorch will be used to simplify implementation of vector operations and embeddings.

\section{Evaluation Plan}

The system will be evaluated using the query and relevance annotations provided in the dataset.

The following evaluation metrics will be used:

\begin{itemize}
    \item Mean Reciprocal Rank at 10 (MRR@10)
    \item Recall@10 and Recall@50
    \item nDCG@10
\end{itemize}

The following systems will be compared:

\begin{enumerate}
    \item BM25 retrieval only,
    \item BM25 retrieval with dense reranking.
\end{enumerate}

This evaluation will help determine whether the semantic reranking stage improves retrieval quality compared to the lexical baseline.

\section{Expected Outcomes}

The expected outcome of the project is a working prototype of a Python code search engine that supports natural language queries. The project will also provide an experimental comparison between lexical-only retrieval and semantic reranking methods.

\section{Timeline}

\begin{itemize}
    \item Week 1--2: dataset preparation and preprocessing
    \item Week 3--4: implementation of BM25 indexing and retrieval
    \item Week 5--6: implementation of dense embeddings and reranking
    \item Week 7: evaluation and experiments
    \item Week 8: analysis of results and project report
\end{itemize}

\begin{thebibliography}{9}

\bibitem{codesearchnet}
Husain, H., Wu, H., Gazit, T., Allamanis, M., and Brockschmidt, M.  
CodeSearchNet Challenge: Evaluating the State of Semantic Code Search.  
\textit{arXiv preprint arXiv:1909.09436}, 2019.  
https://arxiv.org/abs/1909.09436

\bibitem{codebert}
Feng, Z., Guo, D., Tang, D., Duan, N., Feng, X., Gong, M., Shou, L., Qin, B., Liu, T., Jiang, D., and Zhou, M.  
CodeBERT: A Pre-Trained Model for Programming and Natural Languages.  
\textit{Findings of EMNLP}, 2020.  
https://arxiv.org/abs/2002.08155

\bibitem{bm25}
Robertson, S. and Zaragoza, H.  
The Probabilistic Relevance Framework: BM25 and Beyond.  
\textit{Foundations and Trends in Information Retrieval}, 2009.

\bibitem{dense}
Reimers, N. and Gurevych, I.  
Sentence-BERT: Sentence Embeddings using Siamese BERT-Networks.  
\textit{EMNLP}, 2019.

\bibitem{codesearchnetrepo}
CodeSearchNet Dataset Repository.  
https://github.com/github/CodeSearchNet

\end{thebibliography}

\end{document}